\paragraph{Prediction after additional evidence.}
Across 128 matched surface-track cases, four adaptive scalar measurements change combined coordinate RMSE from 0.0578 to 0.0340\,mm. The development-selected fixed mixed schedule gives 0.0423\,mm, camera-only gives 0.0403\,mm, and random selection gives 0.0438\,mm. The paired adaptive-minus-fixed difference is -0.0083\,mm, with 95\% case-bootstrap interval [-0.0110,-0.0058]. Against fixed camera-only, the paired difference is -0.0063\,mm [-0.0083,-0.0043]. The adaptive policy selects 439 camera and 73 probe measurements across 512 decisions. These are complete-cohort comparisons under the declared equal-weight target, not selection of favorable objects.

\paragraph{Different measurements serve different predictions.}
Adaptive selection gives hidden-motion/new-load errors 0.0446/0.0144\,mm. The probe-only policy gives 0.0743/0.0099\,mm; camera-only gives 0.0502/0.0215\,mm. A policy can therefore be useful for changed interactions without being best for reconstructing the original load sequence. The adaptive policy is worse than fixed mixed for new-load and material error; its advantage concerns the combined target. The shared objective makes this tradeoff explicit (Fig.~\ref{fig:acquisition}).

\paragraph{Unknown force magnitude and calibrated evidence.}
The independent 128-case scale-mixture study tests whether apparent material certainty is supplied by the force prior. Under the hierarchical model, four calibrated probes change new-load error from 0.0274 to 0.0141\,mm and material relative error from 0.477 to 0.382. Four views instead give hidden-motion error 0.0626\,mm, compared with 0.0939\,mm initially. Mean log-stiffness/log-force-scale posterior correlation is 0.346 initially and 0.181 after probes. These branches start from the same initial evidence.

After probes, hierarchical force intervals have 88.8\% coverage and 0.158\,N mean width; fixed nominal intervals have 81.9\% coverage and 0.134\,N width. The variance-matched broader Gaussian gives 86.6\% coverage and 0.163\,N width. The hierarchy is thus compared with an explicit wider-prior control, rather than interpreting width alone as calibrated identification.

\paragraph{Dependence, continuous inference and mismatch.}
After adaptive evidence, retaining joint dependence gives mean prediction-interval width 0.0717\,mm versus 0.0863\,mm for the moment-independence control, with coverages 92.0\%/93.9\%. Hidden-motion errors are nearly equal (0.0446/0.0449\,mm); the benefit is interval efficiency, not a large mean-prediction gain. The supplement retains all point and representation controls. On 64 continuous-prior cases, continuous inference gives post-adaptive RMSE 0.0294\,mm versus 0.0350\,mm for fixed mixed measurements. The posthoc paired 95\% interval for their difference is [-0.0088,-0.0027]\,mm. Prediction/force 90\% coverages are 89.7\%/91.2\%. All 192 main fits and 96 seed/particle diagnostics completed; 2 of 144 descriptive sensitivity comparisons exceed the declared thresholds. These checks do not certify convergence. Actions were selected by the finite planner, and the continuous and finite priors differ.

For the finite-dictionary estimator, matched adaptive prediction intervals cover 92.0\% of evaluated coordinates at nominal 90\%. Coverage becomes 86.4\% for continuous parameters, 86.1\% for finer-mesh truth, and 79.2\% for the shorter-correlation force prior. These tests distinguish conditional numerical exactness from robustness to a different generator.

\paragraph{Actual spatial force distributions.}
A separately declared posthoc diagnostic integrates traction errors over the loaded surfaces and all loading times. Joint posterior-mean traction RMSE changes from 10.15 to 7.18\,Pa after adaptive evidence; fixed mixed measurements give 8.04\,Pa. Joint and product-marginal mean fields are nearly identical on this endpoint; the force-field improvement comes from additional evidence. This field metric complements the original one-dimensional load-basis proxy.

\paragraph{A targeted temporal-prior repair.}
On 64 fresh conditional stress cases, mixing correlation length as well as amplitude reduces hidden-motion error after four fixed views from 0.1170 to 0.1007\,mm. Force 90\% coverage rises from 72.4\% to 83.9\%, while width rises from 0.138 to 0.177\,N. The known-covariance oracle gives 91.7\% coverage. Material relative error worsens from 0.476 to 0.506; improved force interpolation does not establish better material recovery. This posthoc follow-up retains the original mismatch failures.

\paragraph{Forward-model discretization.}
In the separate 32-case interior-sensor reference, refining the inverse model changes post-adaptive prediction coverage from 76.7\% to 87.2\%. Overall RMSE changes modestly, 0.0415 to 0.0403\,mm; its paired interval includes zero improvement. Numerical accuracy of the forward model therefore matters for physical uncertainty as well as response error.
